\documentclass[10pt]{article}
\usepackage[margin=0.9in]{geometry}
\usepackage{amsmath,amssymb,amsthm,microtype}
\newtheorem{theorem}{Theorem}
\title{Sampled Coefficient Contraction and Nondivergence}
\author{Collatz proof project}
\date{September 5, 2026}
\begin{document}
\maketitle
\begin{abstract}
We prove in Lean that universal coefficient contraction observed only at
visits to two modulo nine is equivalent to boundedness of every positive
Collatz orbit. The proof selects inverse-coefficient maxima on the sampled
visits and controls intervening coefficients by a factor of sixty-four.
Neither equivalent statement is proved here. Nontrivial cycles are not
excluded, and no mathematical priority claim is made.
\end{abstract}
Let $n_t=T^t(n)$ for the shortcut Collatz map, let $j_t(n)$ count odd
steps, and write $C_t(n)=3^{j_t(n)}/2^t$.
Call $n$ sampled-noncontracting if $n\equiv2\pmod9$ and
$C_t(n)\geq1$ whenever $n_t\equiv2\pmod9$.

\begin{theorem}
Every sampled-noncontracting seed satisfies $C_t(n)\geq1/64$ at all times.
Every positive sampled-noncontracting seed has an unbounded orbit.
\end{theorem}
\begin{proof}
Fix a time $t$ and take the last residue-two visit $k\leq t$; time zero
is such a visit. Our previous residue-height theorem bounds the number
of even steps between $k$ and $t$ by six. The coefficient on that segment
is therefore at least $1/64$. Multiplying by $C_k(n)\geq1$ gives the bound.

Suppose all orbit values were at most $B$. The exact affine identity
$2^t n_t=3^{j_t(n)}n+d_t(n)$ then gives
$d_t(n)/3^{j_t(n)}\leq64B$ for all $t$.
This bounds the ideal correction, which our earlier formal summability
theory proves equivalent to escape for positive seeds. This contradicts
the assumed orbit bound.
\end{proof}

\begin{theorem}
Every unbounded positive orbit has a sampled-noncontracting tail.
\end{theorem}
\begin{proof}
First shift to a visit to two modulo nine, using formal forward coverage.
The shifted orbit is still unbounded. Its inverse coefficient
$f(t)=1/C_t$ tends to zero by the existing summability theorem, and
$f(0)=1$. Consequently $f$ attains a maximum on the subset of times
whose orbit value is two modulo nine: a sufficiently late tail has
$f(t)<1$, so only finitely many earlier sampled times need be compared.
Choose a maximizing sampled time $s$.
For every subsequent sampled visit at time $s+t$, the coefficient cocycle
gives $1/C_t(n_s)=f(s+t)/f(s)\leq1$. Thus $n_s$ is the required seed.
\end{proof}

\begin{theorem}
Every positive orbit is bounded if and only if for every positive
$n\equiv2\pmod9$ some $t$ satisfies $n_t\equiv2\pmod9$ and $C_t(n)<1$.
\end{theorem}
\begin{proof}
Under universal boundedness, failure of sampled contraction contradicts
Theorem 1. Conversely an unbounded orbit supplies, by Theorem 2, a
positive sampled-noncontracting seed, contradicting universal sampled
contraction.
\end{proof}

\paragraph{Scope and remaining work.}
This resolves simultaneous selection of a residue and \emph{sampled}
noncontraction, not selection with full noncontraction at every time.
The latter property cannot be substituted for the former in a proof.
Universal sampled contraction is still an open nondivergence obligation.
The separate candidate that the first sampled contraction forces descent
also remains unproved; the earlier finite scan supplies evidence only.
Neither the equivalence nor its unproved sides exclude nontrivial cycles.

\raggedright
\paragraph{Formal artifacts.}
The module \texttt{Collatz.SampledContraction} contains the definition
\texttt{SampledNeverContracts}, its \texttt{deficit} and \texttt{unbounded}
theorems, \texttt{unbounded\_has\_sampled\_noncontracting\_tail}, and
\texttt{all\_bounded\_iff\_sampled\_contraction}. It imports the project's
proved residue-coverage and noncontracting-tail results. No new axiom or
experimental premise is used.
\end{document}
